%------------------------------------------------------------------------------%
\begin{question}
  Calcule o limite
  \(
    \lim_{(x, y)\to(0, 0)} \frac{x^2-xy}{\,\sqrt{x}-\sqrt{y}\,}
  \)
  se ele existir

  \pause
  \vspace{\baselineskip}
  Assumindo que todos os limites necessários para aplicar as propriedades existam

  (Isso pode não ser verdade)
\end{question}

%------------------------------------------------------------------------------%
\begin{resolution}{Solução}
\begin{align*}
  \onslide<+->{\lim_{(x, y)\to(0, 0)} \frac{x^2-xy}{\,\sqrt{x}-\sqrt{y}\,}}
  \onslide<+->{& = \frac{\displaystyle \lim_{(x, y)\to(0, 0)} \left(x^2-xy\right)}
           {\displaystyle \lim_{(x, y)\to(0, 0)} \left(\sqrt{x}-\sqrt{y}\right)} \\[2mm]}
  \onslide<+->{& = \frac{\displaystyle \lim_{(x, y)\to(0, 0)} x^2 - \lim_{(x, y)\to(0, 0)}xy}
           {\displaystyle \lim_{(x, y)\to(0, 0)} \sqrt{x} - \lim_{(x, y)\to(0, 0)}\sqrt{y}} \\[2mm]}
  \onslide<+->{& =   \frac{0-0}{0-0} }
  \onslide<+->{\;=\; \frac{0}{0} }
  \onslide<+->{\;=\; 1}
\end{align*}

\onslide<+->{\textcolor{red}{\textbf{Não existe divisão por zero !!!}}}

~

~
\end{resolution}

%------------------------------------------------------------------------------%
\begin{resolution}{Solução}
\begin{align*}
  \,\lim_{(x, y)\to(0, 0)} \frac{x^2-xy}{\,\sqrt{x}-\sqrt{y}\,}
  & = \frac{\displaystyle \lim_{(x, y)\to(0, 0)} \left(x^2-xy\right)}
    {\displaystyle \lim_{(x, y)\to(0, 0)} \left(\sqrt{x}-\sqrt{y}\right)} \\[2mm]
  & = \frac{\displaystyle \lim_{(x, y)\to(0, 0)} x^2 - \lim_{(x, y)\to(0, 0)}xy}
    {\displaystyle \lim_{(x, y)\to(0, 0)} \sqrt{x} - \lim_{(x, y)\to(0, 0)}\sqrt{y}} \\[2mm]
  & =  \frac{0-0}{0-0}
  \;=\; \frac{0}{0}
  \qquad\nexists
\end{align*}

\emph{Então o limite não existe?}

~

~
\end{resolution}

%------------------------------------------------------------------------------%
\begin{resolution}{Solução}
  \begin{align*}
    \,\lim_{(x, y)\to(0, 0)} \frac{x^2-xy}{\,\sqrt{x}-\sqrt{y}\,}
    & \color{red}
    = \frac{\displaystyle \lim_{(x, y)\to(0, 0)} \left(x^2-xy\right)}
    {\displaystyle \lim_{(x, y)\to(0, 0)} \left(\sqrt{x}-\sqrt{y}\right)} \\[2mm]
    & \color{red} = \frac{\displaystyle \lim_{(x, y)\to(0, 0)} x^2 - \lim_{(x, y)\to(0, 0)}xy}
    {\displaystyle \lim_{(x, y)\to(0, 0)} \sqrt{x} - \lim_{(x, y)\to(0, 0)}\sqrt{y}} \\[2mm]
    & \color{red} = \frac{0-0}{0-0}
    \;=\; \frac{0}{0}
  \end{align*}

  Não podemos aplicar a propriedade do quociente

  \pause
  \emph{Precisamos remover a indeterminação}

  ~
\end{resolution}

%------------------------------------------------------------------------------%
\begin{resolution}{Solução}
  \onslide<+->{Ao calcular o limite consideramos apenas os pontos dentro do domínio da função}

  \onslide<+->{Para \(x\neq y\)}

  \begin{align*}
    \onslide<+->{\frac{x^2-xy}{\,\sqrt{x}-\sqrt{y}\,}}
    \onslide<+->{& = \frac{\left(x^2-xy\right) \emph{\left(\sqrt{x}+\sqrt{y}\right)}}
                          {\left(\sqrt{x}-\sqrt{y}\right)\emph{\left(\sqrt{x}+\sqrt{y}\right)}} \\[1mm]}
    \onslide<+->{& = \frac{x\left(\emph{x-y}\right)\left(\sqrt{x}+\sqrt{y}\right)}{\emph{x-y}} \\[1mm]}
    \onslide<+->{& = x\left(\sqrt{x}+\sqrt{y}\right)}
  \end{align*}
\end{resolution}

%------------------------------------------------------------------------------%
\begin{resolution}{Solução}
\onslide<+->{Calculando o limite}

\begin{align*}
  \onslide<+->{\lim_{(x, y)\to(0, 0)} \frac{x^2-xy}{\,\sqrt{x}-\sqrt{y}\,}}
  \onslide<+->{& = \lim_{(x, y)\to(0, 0)} x\left(\sqrt{x}+\sqrt{y}\right) \\[1mm]}
  \onslide<+->{& = \left(\lim_{(x, y)\to(0, 0)} x\right)
    \left[
    \left(\lim_{(x, y)\to(0, 0)}\sqrt{x}\right) +
    \left(\lim_{(x, y)\to(0, 0)}\sqrt{y}\right)\right] \\[2mm]}
  \onslide<+->{& = 0\left(\sqrt{0}+\sqrt{0}\right) \\[2mm]}
  \onslide<+->{& = 0}
\end{align*}
\end{resolution}

%------------------------------------------------------------------------------%
